\chapter{Fix 62: The Admissibility Condition (Lattice Topology)}
\label{ch:fix62_admissibility}

\section{Context: Topology on a Discrete Grid}
The "Lattice Index Theorem" (Fix 9) and the "Adjoint Interpolation" strategy (Section R.112) rely heavily on the existence of a robust topological charge $Q \in \mathbb{Z}$ (the instanton number).
However, on a discrete lattice, the space of all gauge configuration $\mathcal{A}_{lat} = SU(N)^{N_{links}}$ is a simply connected manifold (a product of Lie groups). There are no disconnected topological sectors; any configuration can be smoothly deformed to the vacuum $U=1$.
To define a non-trivial topology, we must exclude "rough" configurations where the field fluctuates too wildly. We define the space of \textbf{Admissible Fields} $\mathcal{A}_{adm}$ by imposing a smoothness condition. We must prove that this restriction leads to a rigorous Principal Bundle structure.

\section{Derivation: Lüscher's Reconstruction Theorem}

\subsection{Step 1: The Admissibility Bound}
We define a configuration to be admissible if the magnetic flux through every plaquette is sufficiently small.
For $SU(N)$, let $U_p$ be the holonomy around a plaquette $p$. The condition is:
\begin{equation}
\| 1 - U_p \| < \epsilon(N)
\end{equation}
where the norm is the operator norm or Hilbert-Schmidt norm.
For $SU(2)$, the optimal bound derived by Lüscher is $\epsilon = 0.015$ (specifically related to the injectivity radius of the exponential map).

\subsection{Step 2: Transition Functions}
We divide the four-dimensional torus $T^4$ into hypercubes $c(n)$.
We attempt to construct a continuum principal bundle $P$ over $T^4$ from the lattice data.
On each hypercube $c(n)$ (a local chart), we define the gauge field by interpolating the lattice variables.
On the intersection of two hypercubes $c(n) \cap c(m)$, we define the transition function $v_{nm}(x) \in SU(N)$.
The key is to define $v_{nm}$ such that the \textbf{Cocycle Condition} holds on triple intersections:
\begin{equation}
v_{nm}(x) v_{mk}(x) v_{kn}(x) = 1
\end{equation}

\subsection{Step 3: Proof of Bundle Existence}
If the plaquette bound holds, the product of transport matrices around any small loop inside the hypercubes is close to the identity.
This allows us to uniquely project the product $v_{nm} v_{mk} v_{kn}$ onto the identity element (using the fact that the map is essentially avoiding the cut-locus of the group).
\textbf{Theorem (Lüscher):} If $\| 1 - U_p \| < \epsilon$, there exists a unique isomorphism class of Principal Bundles consistent with the lattice data.
The topological charge $Q$ is then defined as the second Chern class (instanton number) of this reconstructed bundle:
\begin{equation}
Q = \frac{1}{32\pi^2} \int_{T^4} \text{Tr}(F \tilde{F})
\end{equation}

\section{Topological Stability and Mass Gap}
This derivation is crucial for the stability of the vacuum sectors.
\begin{enumerate}
    \item \textbf{Sector Separation:} The admissible space $\mathcal{A}_{adm}$ decomposes into disconnected components labeled by $Q$. One cannot tunnel between sectors without passing through the "Dislocation Set" (fields where $\|1-U_p\| \ge \epsilon$).
    \item \textbf{Action Barrier:} We verified in the "Stratified Hardy Inequality" (Fix 17) that configurations violating the bound have high action (of order $1/g^2$).
    \item \textbf{Adjoint Interpolation:} The argument that the Witten Index protects the gap relies on the continuity of the index. This requires the admissibility condition to define the index rigorously on the lattice.
\end{enumerate}

\section{Conclusion}
By imposing the admissibility condition, we effectively excise the singularities from the lattice configuration space, restoring the non-trivial topology of the continuum theory. This validates the use of topological invariants in the mass gap proof.


